%!TEX root = ../../thesis.tex

\section{Mugiwara}\label{sec:impl-mugiwara}

\code{mugiwara} connects the EfficientZero agent with the MOGI execution environment.
It does this by implementing the \code{Environment} trait for each SUT.
This necessitates the mapping of generic learner actions to executable input and
generated MOGI instrumentation to observations as well as providing an adequate
reward signal.

The project is split into a core library as well as a single binary per target.
The core library includes basic functions common to every target like coverage instrumentation,
a target-agnostic model of a seed corpus, helpers to drive the emulator, diagnostic probes,
and the runner that orchestrates execution of fuzzing campaigns with multiple parameters.
Target binaries only add target-specific functionality.
This includes setup of the emulator with the specified program, executing learner-generated
actions and routing observations and environments from MOGI to the learner.

\subsection{Coverage Instrumentation}\label{sec:impl-coverage}

Coverage is collected by a single block hook that is registered with MOGI.
The hook receives the block address and passes it to a \code{BlockHitTracker},\index{block hit tracker}
which hashes it, folds the hash into one of 4096 buckets and increments that bucket's counter,
similar to AFL++.

The bucket count trades memory against collisions and is target-dependent.
Usually, it is calibrated once manually before starting a larger fuzzing campaign,
since all buffer sizes need to be known beforehand.
For example, cJSON (\cref{sec:target-cjson}) hits 495 unique blocks on startup, so a bucket
size of 512 would lead to almost every block hit colliding with another one, distorting
the observation and derived reward signal.
Coverage numbers are therefore only comparable between runs at the same bucket count.

The tracker counts either blocks or edges.
Block granularity counts basic block hits, while edge granularity counts
transitions between consecutive blocks like \acrshort{afl}, by hashing the
blocks $h_\text{cur}\oplus(h_\text{prev}\gg 1)$, so that for any two basic
blocks $A$ and $B$, $A\rightarrow B$ and $B\rightarrow A$ land in different
buckets and can therefore be differentiated.

Block hits are then converted to a reward according to~\cref{sec:mdp-rewards}.
First, a \code{Coverage} converts raw block hit counts into exponentially spaced levels.
Both terms of the reward introduced in~\cref{sec:mdp-rewards} are then computed from
this score and any additional necessary state.

\subsection{The Seed Corpus}\label{sec:impl-corpus}

The \code{corpus} module implements a target-agnostic seed corpus\index{seed corpus} as described in \cref{sec:mdp-actions}.
As outlined in~\cref{sec:corpus-level-actions}, the goal is a seed corpus with multiple seeds
at a time that can be selected, mutated and retained or discarded.
This is implemented via \code{SeedCorpus}.
Although it is possible for \code{SeedCorpus} to have an arbitrary fixed length,
we evaluate our approach with a single seed only.
In addition to the currently active seeds, the \code{SeedCorpus} also stores the union coverage
of all seeds seen so far as well as the coverage score record.

What the agent practically controls is therefore the position and value of the mutation,
as well as the retain-or-discard decision.
Since actions are represented as categorical variables, a single seed of length $n=16$
has an associated action space of size $2\times 16\times 256=8192$ which is too large
for \gls{mcts} with a few hundred rollout iterations.
We therefore split a single environment step artificially into two steps, where the
action of the first step decides on whether to retain or discard the previous seed as well
as on the mutation position, while the second step decides on the value of the mutation.
The emulator is only run every second step, while the first step returns an empty
observation and plain reward of 0.
This makes the action space size additive and yields a final size of $\max(2\times 16,256)=256$.
For the first step that only has $2\times 16=32$ distinct actions, we reduce the 256 categories
modulo 32.
An episode is fixed at 64 mutations, which equals 128 MDP steps and 64 guest
executions under this model.
The reward is implemented as described in~\cref{sec:mdp-rewards} and derived from the
collected \code{Coverage}.

\subsection{Targets}\label{sec:impl-targets}

Targets are implemented as \code{case}s in \code{mogi\_sysroots}.
Building these targets is done via Dockerfiles and is integrated into the default
\code{Cargo} build process.
Each target is wrapped inside an \code{Environment} that implements the interface
to convert EfficientZero actions to inputs for MOGI and to convert information collected by
MOGI to observations for EfficientZero.

The \code{Environment} also describes how MOGI is set up and what hooks are registered.
Each program is run once until it yields and then a snapshot is created.\index{snapshot}
An environment reset can then be implemented as a simple snapshot restore.

\subsubsection{Input-level action execution}

The execution loop for the input-level action model as introduced in~\cref{sec:input-level-actions}
consists of pushing a single input, usually a byte, into the input stream of the SUT
and then waking it up via MOGI.
The environment is only reset on episode end, i.e., whenever an error condition is
encountered or the step limit of the episode is exceeded.
A single execution step is illustrated in~\cref{lst:mugiwara-step}.

\begin{rust}[caption={One guest execution in the input-action setting},label={lst:mugiwara-step}]
fn execute(&mut self, action: Action) -> (Vec<f32>, bool) {
    push_stdin(&self.setup.emulator, &[action.input]);
    self.block_hit_tracker.borrow_mut().reset();

    let reason = wake_and_run(&mut self.setup.emulator);

    let counts = self.block_hit_tracker.borrow().to_float_vec();
    (counts, reason == ExitReason::Stopped)
}
\end{rust}

\subsubsection{Corpus-level action execution}

During each step, the SUT is re-executed with a given seed.
This is done by resetting the environment and block hit trackers, pushing the
seed data, waking and running the SUT, and collecting any observations and is
illustrated in~\cref{lst:mugiwara-execution}.

\begin{rust}[caption={One guest execution in the corpus-action setting},label={lst:mugiwara-execution}]
fn execute(&mut self, seed: &Seed) -> Coverage {
    restore_snapshot(&mut self.setup.emulator);
    self.block_hit_tracker.borrow_mut().reset();
    push_stdin(&self.setup.emulator, seed);

    wake_and_run(&mut self.setup.emulator);

    let counts = self.block_hit_tracker.borrow().to_float_vec();
    Coverage::from_counts(&counts)
}
\end{rust}

\subsection{Probes}\label{sec:impl-probes}

A probe\index{probe} in this context is a set of hooks (\cref{hooks}) that collect different
kinds of information during execution.
Adding a diagnostic is therefore a matter of attaching a new probe at startup.

Probes are usually backed by a persistent storage, in our case an SQLite database, and contain
names as well as DDLs of the tables they write into, encoded as SQL statements for creating
said tables.
The actual SQLite connection is owned by a background thread which batches multiple
records into transactions and handles shared counters like \code{env\_steps} and \code{train\_batches}
that are included in all tables.
\Cref{tab:impl-probes} lists the probes and what each is for.

\begin{table}[htbp]
  \centering
  \caption[The diagnostic probes]{%
    The diagnostic probes together with the hooks they use and the purpose they serve.
  }
  \label{tab:impl-probes}
  \footnotesize
  \begin{tabularx}{\textwidth}{@{}llX@{}}
    \toprule
    Probe & Hook & Purpose \\
    \midrule
    \code{LossProbe}
      & \code{on\_finished\_training\_batch}
      & the four loss terms of one gradient step \\
    \code{ReturnProbe}
      & \code{on\_finished\_trajectory}
      & return and length of every episode \\
    \code{ActionRateProbe}
      & \code{on\_finished\_trajectory}
      & share of plays landing in a named action set, pooled over the steps where
        that set has a meaning \\
    \code{SeedProbe}
      & \code{on\_finished\_trajectory}
      & the inputs a run produced, recovered from the stored observations \\
    \addlinespace
    \code{CollapseProbe}
      & \code{on\_finished\_trajectory}
      & action played and visit-count entropy per episode-step position --- the
        policy-collapse matrix \\
    \code{ActionIntrospectProbe}
      & \code{on\_finished\_trajectory}
      & the trie of action prefixes self-play actually walked, and how often \\
    \code{PolicyImprovementProbe}
      & \code{on\_search\_roots}
      & divergence of the root visit distribution from the policy prior --- whether
        search improves on the prior or echoes it \\
    \code{ModelIntrospectProbe}
      & \code{on\_finished\_training\_iteration}
      & value, value prefix and prior of every action along a fixed latent rollout,
        per iteration \\
    \addlinespace
    \code{BufferStateProbe}
      & \code{on\_finished\_data\_iteration}
      & replay-buffer occupancy, priority distribution and reward content \\
    \code{SampledBatchProbe}
      & \code{on\_sampled\_batch}
      & what \gls{per} selected: rewarded steps, padding, degeneracy \\
    \code{ReanalyzedBatchProbe}
      & \code{on\_reanalyzed\_batch}
      & the policy and value targets and the priority updates the next gradient step
        sees \\
    \bottomrule
  \end{tabularx}
\end{table}

Three more probes can be installed to dump the episodes as well as pre- and post-reanalyze batches.
They are disabled by default since they produce lots of data in a short amount of time and slow down fuzzing.
To save space in case these probes are required, observations are deduplicated
through a content-addressed store.
There are also two additional tables being written directly from the environment:
the coverage a run has reached, as well as the best execution of each episode.

\subsection{Campaigns}\label{sec:impl-campaigns}

An \code{Evaluation} is the full description of a single run.\index{fuzz campaign!evaluation}
It describes which target is fuzzed, which seed to use, and which settings to apply.
A campaign\index{fuzz campaign} is an ordered queue of evaluations and is drained by a loop while allowing
multiple evaluations to run in parallel, if enough resources are available.

In general, we limit the runs not on wall-clock time, but on the number of data
collection and training iterations.
This makes more sense for our use case since running multiple evaluations in parallel
distorts how long each evaluation was able to run.
Using an iteration bound instead of a wall-clock time bound also allows comparison
across different hardware.

Parallel execution of evaluations is constrained by the hardware, usually the memory.
This can either be too little RAM to hold multiple replay buffers with
the desired capacity, or too little VRAM to store all tensors at once.
Actual GPU compute is also a bottleneck, but only slows down parallel runs and does
not cause them to be killed.
